\documentclass[journal]{IEEEtran}
\usepackage{silence}
\WarningFilter{titlesec}{Non standard sectioning command}
\WarningFilter{caption}{Unknown document class}
\usepackage{amsmath,amssymb}
\usepackage{graphicx}
\usepackage[table]{xcolor}
\usepackage{booktabs}
\usepackage{array}
\usepackage{titlesec}
\usepackage[T1]{fontenc}
\usepackage[utf8]{inputenc}
\usepackage{url}
\usepackage[font=small,labelfont=bf]{caption}
\usepackage{float}
\definecolor{headerblue}{RGB}{220,230,245}
\definecolor{rowgray}{RGB}{248,249,251}

\renewcommand{\thetable}{\arabic{table}}
\renewcommand{\thefigure}{\arabic{figure}}
\renewcommand{\thesection}{\arabic{section}}
\renewcommand{\thesubsection}{\thesection.\arabic{subsection}}

\titleformat{\section}
  {\normalfont\normalsize\bfseries}
  {\thesection.}{0.6em}{}

\titleformat{\subsection}
  {\normalfont\small\bfseries}
  {\thesubsection.}{0.5em}{}

\captionsetup{
    font=footnotesize,
    labelfont=bf,
    textfont=normalfont,
    justification=centering,
    singlelinecheck=true
}


\captionsetup[figure]{
    name=Fig.
}
\raggedbottom
\begin{document}

\title{\LARGE \bfseries
A Simulation-Driven Digital Twin Framework for Decision-Aware Diabetes Modeling Using Causal and Counterfactual Analysis}
\author{
\Large Zarrin Monirzadeh\\[0.6em]
\normalsize Software \& Data Engineer | ML \& AI Systems\\
\normalsize Portland, OR, USA
}
\maketitle

\section*{\centering ABSTRACT}
\noindent
Diabetes mellitus is a complex and heterogeneous metabolic disorder that requires continuous and individualized management. While recent advances in artificial intelligence have improved glucose prediction and risk assessment, most existing approaches remain focused on forecasting rather than supporting actionable decisions.

This paper presents a digital twin framework for multi-type diabetes that integrates temporal modeling, causal reasoning, and counterfactual simulation to enable decision-oriented analysis. The proposed system is designed to conceptually combine clinical variables with temporally structured signals representing physiological and behavioral factors, forming a dynamic patient-specific representation.

The framework maintains a latent state that evolves over time and is used not only to estimate near-term outcomes but also to simulate alternative intervention scenarios. These include variations in dietary intake, physical activity, and treatment adjustments. By modeling intervention effects explicitly, the approach supports interpretable comparison of potential actions rather than relying solely on correlation-based prediction.

To ensure transparency and reproducibility, a proof-of-concept evaluation is conducted using a publicly available diabetes dataset, supplemented by controlled simulated scenarios to illustrate temporal behavior and intervention effects. Experimental results indicate that gradient boosting provides strong predictive performance on structured data, while the simulation component demonstrates how different interventions can influence outcome trajectories.

The primary contribution of this work is a structured and extensible framework that unifies prediction and simulation within a single architecture. Rather than claiming clinical readiness, the study provides a reproducible foundation for future research on decision-support systems in personalized diabetes care.

This work advances the paradigm from predictive artificial intelligence toward decision-aware systems, where models not only forecast outcomes but enable structured evaluation of intervention strategies. By integrating causal reasoning and counterfactual simulation within a unified digital twin architecture, the framework supports actionable decision-making and represents a step toward clinically relevant, patient-specific AI systems.

\vspace{0.6em}
\noindent\textbf{Index Terms:} Digital Twin, Diabetes Mellitus, Causal Inference, Decision Support Systems, Counterfactual Simulation, Time Series Modeling, Personalized Medicine, Clinical Artificial Intelligence.
\vspace{0.8em}
\noindent
\textbf{Preview Version Notice:}
This document provides a concise overview of the proposed framework, focusing on key concepts, methodology, and representative results. 
A comprehensive version including full experimental design, extended results, and implementation details is available upon request.

\vspace{0.3em}
\noindent
\textbf{Contact:} \texttt{monirzadehzarrin@gmail.com}
\section{Introduction}
Diabetes mellitus is a major global health challenge with significant impact on morbidity, mortality, and healthcare systems. It is a heterogeneous group of metabolic disorders rather than a single disease, with substantial variation in underlying physiology and treatment requirements. Type 1 diabetes involves autoimmune destruction of pancreatic beta cells and requires lifelong insulin therapy, while Type 2 diabetes is associated with insulin resistance and progressive metabolic dysfunction. Gestational diabetes occurs during pregnancy and introduces additional maternal and fetal risks.

This limitation reflects a broader gap in current artificial intelligence systems in healthcare. While predictive models provide estimates of future outcomes, they do not directly support decision-making under intervention. In contrast, clinically meaningful intelligence requires the ability to evaluate alternative actions and their potential effects. This distinction motivates a transition from predictive artificial intelligence toward decision-aware systems, in which models are designed not only to forecast outcomes but to support actionable, intervention-driven reasoning.

In practice, diabetes management is inherently decision-driven. Patients and clinicians continuously evaluate interventions such as diet, physical activity, and treatment adjustments under uncertainty. Therefore, intelligent systems must move beyond passive prediction and support actionable decision-making.

The novelty of this work lies in the integration of temporal modeling, causal reasoning, and counterfactual simulation within a unified decision-oriented framework.

The main contributions of this paper are as follows:
\begin{itemize}
\item A modular digital twin architecture adaptable across Type 1, Type 2, and gestational diabetes.
\item A transition from correlation-based prediction to decision-aware modeling, enabling explicit evaluation of intervention strategies through causal and counterfactual reasoning.
\item A reproducible proof-of-concept evaluation using an open dataset and simulated scenarios.
\item A discussion of practical deployment considerations, including interpretability, safety, privacy, and validation.
\end{itemize}

\section{Clinical Computing Foundations}
\subsection{Diabetes as a Multi-Type Control Problem}
From a computational perspective, diabetes management can be formulated as a dynamic control problem under uncertainty. The underlying physiological state evolves over time, observations are incomplete and noisy, external inputs occur frequently, and interventions produce delayed and context-dependent effects.

In Type 1 diabetes, the primary challenge is short-term glucose regulation under the combined influence of insulin administration, food intake, physical activity, stress, and circadian variability. In Type 2 diabetes, longer-term disease progression becomes more prominent, including medication adherence, weight dynamics, and behavioral patterns. Gestational diabetes introduces pregnancy-specific hormonal changes and reduced tolerance for uncertainty due to maternal and fetal safety requirements.


\subsection{Digital Twins in Healthcare}
Medical digital twins are commonly defined as patient-specific computational models that integrate physiological knowledge, individualized data, and continuous state 

However, healthcare digital twins differ fundamentally from industrial systems. The system being modeled is biologically adaptive, partially observable, and strongly influenced by behavioral and contextual factors. The true patient state cannot be directly measured and must instead be inferred. Consequently, digital twins in healthcare must operate under uncertainty, rely on latent state estimation, and incorporate elements of causal structure. These characteristics increase the importance of interpretability, safety, and robustness.

\subsection{Why Causality Matters}
Many current artificial intelligence systems in healthcare learn predictive relationships without explicitly distinguishing causal structure from correlation. In the context of diabetes management, this limitation is particularly important. A model may identify that elevated glucose levels are associated with reduced physical activity, but such associations alone do not provide reliable guidance for intervention.

A digital twin should therefore support explicit intervention analysis. Rather than reporting expected future values alone, it should enable statements such as: given the current patient state, reducing carbohydrate intake from 60 grams to 30 grams is expected to lower peak glucose within a defined time window. This capability distinguishes decision-oriented systems from purely predictive models and represents a central component of the framework proposed in this work.

\section{Related Work}

In addition, prior work by the author on machine learning-based detection of violations in credit card systems highlights the broader relevance of anomaly detection, risk scoring, and decision support under uncertainty~\cite{ref15}. Although the application domain differs, the underlying methodological theme is similar: predictive models become more valuable when they support actionable decisions rather than static classification.

The present work addresses this gap by integrating these research directions into a unified decision-oriented framework. Rather than introducing a fundamentally new learning algorithm, the contribution lies in a structured synthesis of components that are rarely combined within a single system. The proposed approach unifies multi-type diabetes modeling, temporal learning, causal reasoning, counterfactual intervention analysis, and reproducible proof-of-concept evaluation. In this way, the work moves beyond isolated prediction tasks toward a coherent framework for actionable decision support and future clinical translation.

\section{Causal Digital Twin Framework}
\subsection{Design Principles}
The proposed framework is guided by five core principles: patient specificity, multimodality, temporal continuity, counterfactual actionability, and interpretability under safety constraints. Patient specificity emphasizes adaptation to individual physiological characteristics rather than reliance on population-level averages. Multimodality reflects the integration of clinical, physiological, and behavioral signals. Temporal continuity captures longitudinal dependencies instead of isolated observations. Counterfactual actionability enables evaluation of alternative interventions, while interpretability under safety constraints ensures that recommendations remain explainable and clinically appropriate.

\subsection{System Overview}
The proposed framework follows a modular pipeline consisting of data ingestion, state representation, predictive modeling, and counterfactual simulation. Clinical, physiological, and behavioral inputs are first harmonized and transformed into structured representations. A latent state representation is then constructed to capture patient-specific dynamics, which are subsequently used for prediction and simulation of intervention outcomes.

\begin{figure}[htbp]
\centering
\includegraphics[width=0.95\linewidth]{Fig2.png}
\caption{Proposed digital twin architecture illustrating the flow from data ingestion and preprocessing to latent state representation, predictive modeling, and counterfactual simulation for intervention analysis.}
\label{fig:architecture}
\end{figure}

\subsection{Data Layer}
The input layer integrates heterogeneous variables across three categories: clinical, physiological, and behavioral. Clinical variables include demographic information, diagnosis type, medication regimen, and comorbidities. Physiological variables include glucose measurements and, when available, additional biometric signals. Behavioral variables include meal intake, physical activity, sleep patterns, and other contextual factors.

In the current implementation, multimodal integration is demonstrated conceptually, while experimental validation is conducted on structured benchmark data.

Due to differences in sampling rates, noise levels, and data reliability, preprocessing involves temporal alignment, normalization, missing-value imputation, and feature fusion. For real-time applications, a rolling time window is employed to capture recent dynamics and maintain temporal consistency.

\begin{figure}[htbp]
\centering
\includegraphics[width=0.9\linewidth]{Fig3.png}
\caption{Mapping between input features and predictive models, including both classical machine learning methods and neural network-based approaches within the digital twin framework.}
\label{fig:model_mapping}
\end{figure}

It is important to emphasize that these values should be interpreted as qualitative demonstrations of intervention effects rather than quantitative clinical estimates. The results are derived from a simulated illustrative setup and do not represent validated clinical outcomes. Their purpose is to demonstrate how predictive modeling can be extended into actionable decision support through counterfactual reasoning.

Overall, these findings support the feasibility of the proposed architecture while highlighting the need for future evaluation on longitudinal clinical datasets.

All experimental outputs, including generated trajectories and tabular results (CSV/XML formats), are publicly available in the accompanying repository~\cite{ref41}.

\section{Discussion}
This paper presented a personalized digital twin framework for multi-type diabetes that advances artificial intelligence from predictive modeling toward decision-aware clinical support. Rather than estimating outcomes in isolation, the framework enables structured evaluation of intervention strategies through counterfactual simulation and causal reasoning. This approach provides a foundation for actionable, patient-specific decision-making in healthcare systems.

The experimental results support this transition at a conceptual level. The benchmark evaluation demonstrates that established machine learning models can provide stable predictive baselines on structured data, while the counterfactual simulation illustrates how these predictions can be extended to intervention analysis. Together, these components highlight the feasibility of integrating prediction and decision support within a unified framework.

From a broader perspective, this work reflects a conceptual shift in healthcare artificial intelligence from passive prediction toward active decision support. By enabling structured comparison of alternative interventions, the proposed framework aligns with emerging needs in precision medicine, where patient-specific decision-making is essential. This shift has implications not only for diabetes management but also for a wider class of chronic diseases that require continuous monitoring and adaptive intervention strategies.

\section{Limitations}
This study has several limitations. First, the benchmark dataset does not include continuous glucose monitoring time-series data and therefore does not support real-time glycemic forecasting. As a result, the temporal and intervention-related components are evaluated using simulated scenarios rather than real longitudinal observations.

Second, the counterfactual trajectories are generated using simplified parametric models and are intended for illustrative purposes only. While these simulations follow physiologically plausible patterns, they do not capture the full complexity of real patient responses.

Third, the causal structure is presented as a conceptual design rather than a fully estimated causal model derived from observational or interventional data. This limits the ability to draw causal conclusions in a strict statistical sense.

The objective of this study is not to establish clinical performance, but to provide a transparent and reproducible foundation for simulation-driven digital twin research in diabetes.

\section{Future Work}
Future work should evaluate the framework using longitudinal continuous glucose monitoring datasets with integrated annotations for meals, physical activity, sleep, medication, and insulin. Such data would enable direct assessment of short-term forecasting, patient-state adaptation, and intervention simulation under realistic conditions.


\section{Conclusion}
This paper presented a personalized digital twin framework for multi-type diabetes that emphasizes simulation-driven decision support rather than prediction alone. The framework integrates clinical, physiological, and behavioral data, maintains a patient-specific latent state, and applies counterfactual simulation to compare feasible interventions prior to real-world execution.
--------

This shift toward decision-aware artificial intelligence enables a new class of clinical systems that move beyond passive monitoring toward proactive, personalized intervention planning.

\begin{thebibliography}{99}

\bibitem{ref15}
Z. Monirzadeh, “Detection of Violations in Credit Cards of Banks Using Machine Learning Techniques,” 2018.

\bibitem{ref41}
Z. Monirzadeh, “A Simulation-Driven Digital Twin Framework for Diabetes Modeling: Code and Experiments,”
GitHub, 2026. [Online]. Available:
\url{https://github.com/zarrinmonirzadeh/diabetes-ai-ml-digital-twin}

\end{thebibliography}

\vspace{2em}

\begin{center}
{\Large \bfseries Full Article Available Upon Request}
\end{center}

\vspace{0.8em}

\begin{center}
{\normalsize
For access to the complete manuscript, including full methodology, extended experiments, and implementation details, please contact:}
\end{center}

\vspace{0.5em}

\begin{center}
{\large \ttfamily monirzadehzarrin@gmail.com}
\end{center}

\vspace{1em}
\end{document}